\section{Experiments}
\label{sec:experiments}

The experiments examine whether model refinement and command
correction bring the physical cable motion closer to a fixed
simulated reference. One target and one initial whipping plan
are used throughout, so changes in performance are measured
against the same desired trajectory and timing.

\subsection{Setup and Refinement Protocol}
\label{subsec:experimental_setup}

The platform comprises a 145\,g quadrotor and a 17\,g cable
assembly, including tracking markers. The cable length is
0.9525\,m. OptiTrack records vehicle pose and cable-marker
positions, and controller logs record the executed commands.
The tracked vehicle origin starts at $(0,0,1.4)$\,m, with
$\mathbf p_{\mathrm{tar}}=(1.25,0,1.25)^\top$\,m and
$\mathbf e=(1,0,0)^\top$. Commands are supplied at 30\,Hz;
cable measurements do not modify them during execution.

The predictor uses $N=12$ cable nodes and produces rollout
samples at 150\,Hz. The command spline has 12 control points
over $T=1.5$\,s. Correction tracks the original reference through
$T_h=34/30$\,s, with weights
$(w_{\mathrm{tip}},w_q,w_d)=(1,0.1,0.01)$
in~\eqref{eq:reference_tracking_cost}. Each exported sequence
also includes braking, return to launch, and a final hold.

Five executions were collected for each of three successive
model--command pairs, denoted M0, M1, and M2. M0 uses the
preliminary model and initial plan. M0 measurements refine
$\mathcal M_1$, which corrects the original commands. M1
measurements then refine $\mathcal M_2$, which corrects the
M1 commands toward the same M0 reference. No new MPPI search
or change of target is introduced between batches.

In each of the M0 and M1 batches, recordings 1, 2, and 4
contribute to fitting alongside earlier data. Recordings 3
and 5 support development checks; the two M1 recordings also
support M2 model selection. They are therefore not independent
test data. The selected M2 retains the M1 vehicle network,
updates the nominal vehicle response and cable model, and
is frozen before the five M2 flights. No M2 recording is
used for fitting or command correction. All 15 executions
remain in the primary physical comparison.

\subsection{Evaluation Metrics}
\label{subsec:experimental_metrics}

Physical tracking is evaluated against the original predictions
$\mathbf r_{\mathrm{tip}}^0(t)$ and $\mathbf p_q^0(t)$, rather
than predictions regenerated after each model update.
For $n$ valid measurement times $t_j\in[0,T_h]$, tip-reference
error is
\[
E_{\mathrm{tip}}=
\sqrt{\frac{1}{n}\sum_{j=1}^{n}
\|\mathbf r_{\mathrm{tip}}^{\mathrm{meas}}(t_j)
-\mathbf r_{\mathrm{tip}}^0(t_j)\|_2^2}.
\]
The reference is interpolated at these timestamps.
Replacing the tip positions by measured and reference vehicle
positions defines $E_q$. Both are three-dimensional RMSEs.
Clock offsets are estimated from measured vehicle streams;
trajectories are not shifted to improve target agreement or
spatially aligned to the reference.

Target approach is measured by $d_{\min}$ from
\eqref{eq:target_distance} over $[0,1.5]$\,s and by
$d_\star=\|\mathbf r_{\mathrm{tip}}^{\mathrm{meas}}(t_\star)
-\mathbf p_{\mathrm{tar}}\|_2$ at the original planned strike
time $t_\star=1.11725$\,s. These distinguish proximity from
arrival at the planned time. Tip speed at $t_\star$ is estimated
from local derivatives of measured positions. No binary hit
threshold is imposed. Missing observations are masked, and
statistics are computed per execution before taking the
batch mean and sample standard deviation.

\subsection{Physical Reference Tracking and Target Approach}
\label{subsec:physical_results}

\begin{table}[t]
\centering
\caption{Physical performance: mean $\pm$ sample standard
deviation over five executions per batch. Distances are in cm;
tip speed is in m/s. All batches track the same M0 reference.}
\label{tab:physical_performance}
\small
\setlength{\tabcolsep}{3pt}
\begin{tabular}{@{}lccc@{}}
\toprule
Metric & M0 & M1 & M2 \\
\midrule
$d_\star$ & $24.12\pm6.94$ & $11.18\pm4.04$ & $13.10\pm7.52$ \\
$d_{\min}$ & $14.73\pm6.02$ & $7.83\pm1.82$ & $8.60\pm6.55$ \\
$E_{\mathrm{tip}}$ & $18.09\pm1.80$ & $16.06\pm2.34$ & $13.28\pm2.60$ \\
$E_q$ & $15.66\pm3.58$ & $11.17\pm2.87$ & $12.63\pm5.65$ \\
Tip speed & $5.98\pm0.53$ & $5.85\pm0.32$ & $5.82\pm0.79$ \\
\bottomrule
\end{tabular}
\end{table}

Table~\ref{tab:physical_performance} shows that the first
refinement substantially improves target approach: mean
$d_\star$ decreases from 24.12 to 11.18\,cm, and mean
$d_{\min}$ decreases from 14.73 to 7.83\,cm. Tip-reference
RMSE decreases more modestly, from 18.09 to 16.06\,cm.
The second refinement reduces it further to 13.28\,cm,
a 17.4\% reduction relative to M1. Mean tip speed remains
similar across batches.

Target accuracy does not improve monotonically. M2 has larger
mean $d_\star$ and $d_{\min}$ than M1, and the sample standard
deviation of $d_\star$ increases from 4.04 to 7.52\,cm.
Thus, improved average trajectory tracking does not ensure
a more accurate or repeatable target approach. This distinction
also reflects the correction objective, which penalizes
trajectory error over an interval rather than target error alone.

\noindent\emph{Exploratory subset.}
After reviewing the outcomes, M2 recordings 2 and 4 were
selected for an additional analysis. Their mean $d_\star$,
$d_{\min}$, and $E_{\mathrm{tip}}$ are 6.17, 4.41, and
10.57\,cm, respectively; mean tip speed is 5.84\,m/s.
Hardware defects in recordings 1, 3, and 5 were suspected
but not independently established. These two selected
executions therefore do not replace the five-trial result
or establish typical M2 performance.

\subsection{Model Prediction and Limitations}
\label{subsec:prediction_results}

On the two M1 development recordings used for model selection,
the selected M2 reduces mean vehicle-prediction RMSE from
7.46 to 6.72\,cm and command-driven tip-prediction RMSE from
9.95 to 8.46\,cm relative to M1. These comparisons use the
same recorded commands, observation masks, and initial states
estimated from preceding measurements. They diagnose model
selection; they are neither independent test results nor
tracking errors relative to the planned motion.

The physical study comprises successive batches at one target,
not randomized paired trials. Planning assumes a stationary
hanging cable, whereas small tip motions were observed by the
operator during real hover. The resulting initial-state
mismatch may affect repeatability. Propeller airflow is a
possible explanation for this motion, but its contribution
to the measured errors has not been isolated.
Timing uncertainty also matters: for M2 recording 2,
perturbing clock alignment by $\pm10$\,ms changes $d_\star$
from 2.64\,cm to 6.80--7.15\,cm. The results support improved
tracking of this fixed reference, while broader robustness
and the sources of trial-to-trial variability remain open.

